%%% No Finite Local Potential Function Suffices for the Accelerated
%%% Collatz Map: An Impossibility Result

\documentclass[11pt,a4paper]{article}

%% ---------- packages ----------
\usepackage[utf8]{inputenc}
\usepackage{amsmath,amssymb,amsthm,mathtools}
\usepackage[margin=2.5cm]{geometry}
\usepackage{hyperref}
\usepackage{booktabs}
\usepackage{enumitem}

%% ---------- theorem environments ----------
\newtheorem{theorem}{Theorem}[section]
\newtheorem{proposition}[theorem]{Proposition}
\newtheorem{lemma}[theorem]{Lemma}
\newtheorem{corollary}[theorem]{Corollary}
\theoremstyle{definition}
\newtheorem{definition}[theorem]{Definition}
\newtheorem{example}[theorem]{Example}
\theoremstyle{remark}
\newtheorem{remark}[theorem]{Remark}

%% ---------- macros ----------
\DeclareMathOperator{\val}{v}
\newcommand{\N}{\mathbb{N}}
\newcommand{\Z}{\mathbb{Z}}
\newcommand{\R}{\mathbb{R}}
\newcommand{\Nodd}{\N_{\mathrm{odd}}}

%% ---------- metadata ----------
\title{\Large No Finite Local Potential Function Is a Strict Lyapunov Function\\
for the Accelerated Collatz Map:\\[4pt]
\large An Impossibility Result}
\author{Lightman Chang\\
Independent Researcher\\
\texttt{lightman.chang@gmail.com}}
\date{}

%% =================================================================
\begin{document}
\maketitle

%% =================================================================
%% ABSTRACT
%% =================================================================
\begin{abstract}
A standard strategy for attacking the Collatz conjecture is to seek a
\emph{Lyapunov function}: a real-valued map on the positive integers
that strictly decreases under the accelerated Collatz dynamics
$S(n) = (3n+1)/2^{\val_2(3n+1)}$.  Any such function would force the
absence of nontrivial cycles and divergent orbits.  A natural family of
candidates consists of the \emph{finite local} potential functions of
the form $V(n) = \log_2 n + g(n \bmod 2^k)$, where $k \geq 1$ is a
fixed precision and $g \colon \Z / 2^k \Z \to \R$ is an arbitrary
correction term.  We prove (Theorem~\ref{thm:main}) that for every
fixed $k \geq 1$ and every choice of $g$, there exist infinitely many
odd integers $n$ on which $V$ does not decrease under~$S$, so $V$ fails
to be a strict Lyapunov function for $S$.  The obstruction is structural:
the trailing-$1$ block can have arbitrary length, and a finite-precision
correction can compensate at most $O(1)$ growth, while the trailing-$1$
block of length $\ell$ produces a multiplicative growth of
$(3/2)^\ell$ before any cancellation.  Our result rules out an entire
class of potential-function attacks on Collatz and complements
Conway's~\cite{Conway1972} undecidability theorems for generalized
Collatz problems.
\end{abstract}

\noindent\textbf{Keywords:}~Collatz map, $3x{+}1$ problem, Lyapunov
function, impossibility result, $2$-adic valuation.

\smallskip
\noindent\textbf{MSC 2020:}~Primary 11B37; Secondary 37B99, 11A63.

%% =================================================================
\section{Introduction}\label{sec:intro}
%% =================================================================

The Collatz conjecture asserts that for the map
$T \colon \N \to \N$ defined by $T(n) = n/2$ for $n$ even and
$T(n) = 3n + 1$ for $n$ odd, every orbit reaches~$1$.  After nearly a
century, it remains open; the strongest known density result is due to
Tao~\cite{Tao2022}, and the standard surveys are
Lagarias~\cite{Lagarias1985,Lagarias2010}.  Conway~\cite{Conway1972}
showed that the natural family of generalizations of $T$ is
universal (Turing-complete), so any proof of Collatz must exploit
specific arithmetic features of the constants $(3, 2, 1)$.

A frequently proposed strategy is to construct a real-valued function
$V \colon \Nodd \to \R$ that strictly decreases under the accelerated
Collatz map
\begin{equation}\label{eq:S-def}
  S \colon \Nodd \to \Nodd, \qquad
  S(n) \;=\; \frac{3n + 1}{2^{\val_2(3n + 1)}}.
\end{equation}
Such a $V$ is called a (strict) \emph{Lyapunov function} for $S$.
The existence of any such $V$ that is bounded below (or merely
well-founded) would immediately rule out cycles and divergent orbits,
and hence prove the Collatz conjecture.

\paragraph{The natural ansatz.}
On any orbit, $S$ scales the iterate by approximately $3 / 2^{\val_2}$,
so the leading-order behavior of $\log_2 n$ along an orbit is governed
by $\log_2 3 - \val_2 \approx 1.585 - \val_2$.  This suggests an
ansatz of the form
\begin{equation}\label{eq:V-form}
  V(n) \;=\; \log_2 n \;+\; g(n \bmod 2^k),
\end{equation}
where $k \geq 1$ is a fixed integer (the ``locality precision'') and
$g \colon \Z / 2^k \Z \to \R$ is a correction term whose role is to
absorb the fluctuations of $\val_2(3n + 1)$ that depend on the low
bits of $n$.  Many variants of this idea are folkloric; small-$k$
versions can be checked by hand and seen to fail.

\paragraph{What we prove.}
We prove that the failure is structural: for \emph{every} fixed $k$ and
\emph{every} choice of $g$, the function $V$ in~\eqref{eq:V-form} fails
to be a strict Lyapunov function for $S$.  More precisely, there exist
infinitely many odd integers $n$ on which $V(S(n)) \geq V(n)$.  The
obstruction comes from the trailing-$1$ blocks: an odd integer of the
form $n = 2^\ell - 1$ has a deterministic prefix of $\ell - 1$
``bad'' steps, each contributing $\log_2(3/2)$ of growth before any
cancellation, while $g$ takes values in a fixed finite set and can
contribute at most a fixed constant of cancellation.

\paragraph{What we do not prove.}
We do not prove (and do not believe) that no Lyapunov function exists
at all.  The result rules out only the natural finite-precision local
ansatz~\eqref{eq:V-form}.  It does not rule out (a) potentials of the
form $\log_2 n + g(n \bmod 2^{k(n)})$ with $k(n) \to \infty$, (b)
potentials taking $g$ to be a function of more elaborate features
(e.g., $\val_2(n+1)$ itself, or block-structure of the binary
expansion), or (c) ordinal-valued potentials.  We discuss these
extensions in Section~\ref{sec:disc}.

\paragraph{Why the result is useful.}
Various authors have tried and abandoned the
ansatz~\eqref{eq:V-form} for small $k$, often after lengthy ad-hoc case
analysis.  Our theorem replaces this by a one-sentence obstruction:
the family is dead because $\log_2(3/2) > 0$ and trailing-$1$ blocks
are unbounded in length.  This frees the search for
Lyapunov-style attacks on Collatz to focus on genuinely
non-local potentials.

\paragraph{Paper structure.}
Section~\ref{sec:prelim} fixes notation and establishes the closed-form
behavior of $S$ on the trailing-$1$ block (the only Collatz-specific
input we need).  Section~\ref{sec:main} proves
Theorem~\ref{thm:main}.  Section~\ref{sec:disc} discusses extensions and
limitations.

%% =================================================================
\section{Preliminaries}\label{sec:prelim}
%% =================================================================

We write $\Nodd = \{1, 3, 5, \dots\}$ for the positive odd integers and
$\val_2(m)$ for the $2$-adic valuation of $m \in \Z \setminus \{0\}$.
For an odd integer $n \geq 1$ we set
\[
  t(n) \;=\; \val_2(n + 1) \;\geq\; 1,
\]
the number of trailing $1$-bits in the binary expansion of $n$.  Thus
$n = 2^{t(n)} m' - 1$ with $m'$ odd.

\begin{definition}[Strict Lyapunov function]\label{def:lyap}
A function $V \colon \Nodd \to \R$ is a \emph{strict Lyapunov function}
for $S$ if $V(S(n)) < V(n)$ for every $n \in \Nodd$ with $n > 1$.
\end{definition}

\begin{definition}[Finite local potential]\label{def:loc-pot}
For a fixed integer $k \geq 1$ and a function
$g \colon \Z / 2^k \Z \to \R$, the \emph{finite local potential of
precision $k$ with correction $g$} is
\[
  V_{k, g}(n) \;=\; \log_2 n \;+\; g\bigl(n \bmod 2^k\bigr),
  \qquad n \in \Nodd.
\]
\end{definition}

We will need only one Collatz-specific fact, recorded here as a lemma;
it is a special case of a folkloric closed form (see, e.g.,
\cite[Lemma~2.1]{Lagarias1985}).

\begin{lemma}[Trailing-$1$ block]\label{lem:trailing}
Let $n \geq 3$ be odd with $t(n) = t \geq 2$.  Then for $j = 1, 2,
\dots, t-1$ the iterate $S^j(n)$ is well-defined as an odd integer and
\[
  S^j(n) \;=\; 2^{t - j} \cdot 3^j \cdot m' \;-\; 1,
\]
where $n = 2^t m' - 1$.  In particular,
\begin{equation}\label{eq:Sj-growth}
  \log_2 S^j(n) \;=\; \log_2 n \;+\; j \cdot \log_2(3/2)
  \;+\; \varepsilon_j(n),
\end{equation}
where $|\varepsilon_j(n)| \leq 2^{-(t - j) - \log_2 m'}$ tends to zero
as $n \to \infty$ along $n = 2^t m' - 1$ with $t - j \geq 1$ and $m'$
fixed.
\end{lemma}

\begin{proof}
Write $n = 2^t m' - 1$.  Then $3n + 1 = 3 \cdot 2^t m' - 2 = 2(3 \cdot
2^{t-1} m' - 1)$, and since $3 \cdot 2^{t-1} m' - 1$ is odd (as
$2^{t-1} m'$ is even when $t \geq 2$), we have $\val_2(3n+1) = 1$ and
$S(n) = 3 \cdot 2^{t-1} m' - 1 = 2^{t-1} \cdot (3 m') - 1$.  But
$3 m'$ is odd (product of two odd numbers), so $S(n)$ has the same
form $S(n) = 2^{t - 1} m_1' - 1$ with $m_1' = 3 m'$ odd and
$t(S(n)) = t - 1$.  Iterating,
$S^j(n) = 2^{t - j} \cdot 3^j m' - 1$ for $0 \leq j \leq t - 1$, valid
because $t - j \geq 1$ keeps $S^j(n)$ in the trailing-$1$ regime.

For~\eqref{eq:Sj-growth},
\[
  S^j(n) = 2^{t-j} \cdot 3^j m' - 1
  = 3^j \cdot \frac{2^t m' - 2^j}{2^j}
  = \frac{3^j}{2^j} (n + 1) - 1
  = \frac{3^j}{2^j} n + \frac{3^j - 2^j}{2^j},
\]
so
\[
  \log_2 S^j(n) \;=\; \log_2 n + j \log_2(3/2)
  + \log_2\!\left(1 + \frac{3^j - 2^j}{2^j n}\right).
\]
For $n \geq 2^t \cdot 1 - 1$ and $j \leq t - 1$ the third term is
$O(2^{-(t-j)})$ uniformly in $j$, giving the claimed bound on
$\varepsilon_j(n)$.
\end{proof}

\begin{remark}
We use~\eqref{eq:Sj-growth} only to control the dominant
$j \log_2(3/2)$ term; the error $\varepsilon_j(n)$ will be made
arbitrarily small by taking $n$ large.  No effort to optimize the
constant in the error bound is needed.
\end{remark}

%% =================================================================
\section{The Impossibility Theorem}\label{sec:main}
%% =================================================================

\subsection{Statement}

\begin{theorem}[Main result]\label{thm:main}
Let $k \geq 1$ be a fixed integer.  For every function
$g \colon \Z / 2^k \Z \to \R$, the finite local potential
$V_{k, g}(n) = \log_2 n + g(n \bmod 2^k)$ fails to be a strict
Lyapunov function for the accelerated Collatz map $S$.  More precisely,
there exist infinitely many odd integers $n > 1$ such that
\[
  V_{k, g}\bigl(S(n)\bigr) \;\geq\; V_{k, g}(n).
\]
\end{theorem}

\subsection{Proof}

\begin{proof}[Proof of Theorem~\ref{thm:main}]
Fix $k \geq 1$ and an arbitrary function $g \colon \Z / 2^k \Z \to \R$.
Since $\Z / 2^k \Z$ is finite, the function $g$ is bounded; set
\begin{equation}\label{eq:def-G}
  G \;:=\; \max_{r \in \Z / 2^k \Z} g(r) \;-\; \min_{r \in \Z / 2^k \Z}
  g(r) \;\in\; [0, \infty).
\end{equation}
Note that $G$ depends on $g$ but is finite (since $g$ takes only $2^k$
values).

Choose any integer $\ell$ large enough that
\begin{equation}\label{eq:ell-choice}
  \ell \cdot \log_2(3/2) \;>\; G \;+\; 1.
\end{equation}
This is possible because $\log_2(3/2) = \log_2 3 - 1 \approx 0.585$ is
strictly positive, so the left-hand side grows without bound in $\ell$.
Concretely, $\ell \geq \lceil (G + 1) / \log_2(3/2) \rceil$ suffices.

For each integer $M \geq 1$ define the odd integer
\begin{equation}\label{eq:def-n}
  n_M \;:=\; 2^{\ell + k} \cdot M \;-\; 1.
\end{equation}
Then $n_M$ has the following properties.

\medskip
\noindent\textbf{(i)} $n_M$ is odd, and $t(n_M) = \ell + k \geq \ell$.
Indeed, $n_M + 1 = 2^{\ell + k} M$, so $\val_2(n_M + 1) = \ell + k$
because $M$ may have its own factors of $2$, but we only used
$\ell + k$ explicitly; if $M$ is itself even the trailing-$1$ count
would be larger, which only helps.  Restrict to odd $M$ to ensure
$t(n_M) = \ell + k$ exactly; we will make this restriction below.

\medskip
\noindent\textbf{(ii)} The residue $n_M \bmod 2^k$ is independent of
$M$: indeed, $n_M \equiv -1 \equiv 2^k - 1 \pmod{2^k}$ for every $M$,
since $2^{\ell + k} M \equiv 0 \pmod{2^k}$.  Consequently,
\[
  g(n_M \bmod 2^k) \;=\; g(2^k - 1)
  \quad\text{is the same constant for every } M.
\]

\medskip
\noindent\textbf{(iii)} By Lemma~\ref{lem:trailing} applied with
$t = \ell + k$ and $j = \ell$, the iterate $S^\ell(n_M)$ is well-defined
and satisfies
\[
  S^\ell(n_M) \;=\; 2^{k} \cdot 3^\ell \cdot M \;-\; 1.
\]
This number is odd, has $t(S^\ell(n_M)) = k$ trailing $1$-bits when
$M$ is odd (since $3^\ell M$ is then odd), and most importantly
\[
  S^\ell(n_M) \;\equiv\; 2^k - 1 \;\equiv\; -1 \pmod{2^k},
\]
so $g(S^\ell(n_M) \bmod 2^k) = g(2^k - 1)$ as well.

\medskip
Restricting now to odd $M \geq 1$, we have for every such $M$:
\begin{align}
  V_{k, g}(S^\ell(n_M)) - V_{k, g}(n_M)
  &= \bigl[\log_2 S^\ell(n_M) - \log_2 n_M\bigr]
   + \bigl[g(S^\ell(n_M) \bmod 2^k) - g(n_M \bmod 2^k)\bigr]
   \notag \\
  &= \bigl[\log_2 S^\ell(n_M) - \log_2 n_M\bigr]
   + 0,
   \label{eq:diff-V}
\end{align}
where the second bracket vanishes by~(ii) and~(iii).

By Lemma~\ref{lem:trailing} (specifically~\eqref{eq:Sj-growth}),
\begin{equation}\label{eq:diff-log}
  \log_2 S^\ell(n_M) - \log_2 n_M
  \;=\; \ell \cdot \log_2(3/2) + \varepsilon(M),
\end{equation}
where $\varepsilon(M) \to 0$ as $M \to \infty$ (the error
$\varepsilon_\ell(n_M)$ from Lemma~\ref{lem:trailing} satisfies
$|\varepsilon(M)| \leq C / (2^k M)$ for a constant $C$ depending on
$\ell$, since $n_M \asymp 2^{\ell + k} M$).  Choose $M_0$ so large that
$|\varepsilon(M)| < 1$ for all $M \geq M_0$.

Combining \eqref{eq:diff-V}, \eqref{eq:diff-log} and the
choice~\eqref{eq:ell-choice}, we obtain that for every odd $M \geq M_0$,
\begin{equation}\label{eq:V-grows}
  V_{k, g}(S^\ell(n_M)) - V_{k, g}(n_M)
  \;=\; \ell \cdot \log_2(3/2) + \varepsilon(M)
  \;>\; G + 1 - 1 \;=\; G \;\geq\; 0.
\end{equation}

So far we have shown that the \emph{$\ell$-step} potential difference
$V_{k, g}(S^\ell(n_M)) - V_{k, g}(n_M)$ is strictly positive.  We now
deduce the existence of a single step on which $V_{k, g}$ does not
decrease.  Telescoping,
\[
  V_{k, g}(S^\ell(n_M)) - V_{k, g}(n_M)
  \;=\; \sum_{j = 0}^{\ell - 1}
    \bigl[V_{k, g}(S^{j+1}(n_M)) - V_{k, g}(S^j(n_M))\bigr].
\]
Since the left-hand side is strictly positive by~\eqref{eq:V-grows},
at least one summand on the right-hand side is strictly positive: there
exists an index $j^*(M) \in \{0, 1, \dots, \ell - 1\}$ for which
\[
  V_{k, g}\bigl(S^{j^*(M) + 1}(n_M)\bigr) \;>\;
  V_{k, g}\bigl(S^{j^*(M)}(n_M)\bigr).
\]
In particular, on the odd integer $\tilde n_M := S^{j^*(M)}(n_M)$ the
potential strictly fails to decrease.

It remains to verify that the values $\{\tilde n_M : M \in \Nodd,
M \geq M_0\}$ are infinite (i.e., that we have produced
\emph{infinitely many} witnesses).  Since
$\log_2 \tilde n_M \geq \log_2 n_M - O(\ell) \to \infty$ as
$M \to \infty$ (the bad-step segment of length at most $\ell - 1$ can
shrink the iterate by at most a factor of $(2/3)^{\ell - 1}$ in
log-scale, a finite amount), the set $\{\tilde n_M\}$ is unbounded and
hence infinite.

This proves the existence of infinitely many odd integers $n$ on which
$V_{k, g}(S(n)) \geq V_{k, g}(n)$, as claimed.
\end{proof}

\subsection{Quantitative consequence}

\begin{corollary}[Required precision grows]\label{cor:precision}
Fix any constant $A > 0$ and let $k \geq 1$.  For every $g \colon
\Z / 2^k \Z \to \R$ with $\max g - \min g \leq A$, the function
$V_{k, g}$ admits infinitely many odd integers $n$ on which
$V_{k, g}(S(n)) - V_{k, g}(n) \geq 0$.  In particular, no uniform bound
on the oscillation of $g$ across all $k$ produces a Lyapunov function;
the only escape from Theorem~\ref{thm:main} is to allow $g$'s
oscillation $G(g)$ to grow without bound in~$k$.
\end{corollary}

\begin{proof}
Direct from~\eqref{eq:ell-choice}: any uniform bound $G \leq A$ allows
us to choose $\ell$ once for all $k$, with the construction of
Theorem~\ref{thm:main} going through as written.
\end{proof}

%% =================================================================
\section{Discussion}\label{sec:disc}
%% =================================================================

\paragraph{What is ruled out.}
Theorem~\ref{thm:main} rules out the entire family of Lyapunov
candidates of the form $V_{k, g} = \log_2 n + g(n \bmod 2^k)$ with
$k$ fixed and $g$ a function of the residue class.  This is, to our
knowledge, the most natural class of ``elementary'' Lyapunov
candidates that has been considered in the literature.

\paragraph{What is not ruled out.}
Our result leaves several variants open.

\begin{enumerate}[label=(\arabic*),leftmargin=2em]
  \item \emph{Adaptive precision.}  A potential of the form
    $\log_2 n + g_n(n \bmod 2^{k(n)})$ with $k(n) \to \infty$ as
    $n \to \infty$, and with $g_n$ permitted to vary, is not covered.
    Corollary~\ref{cor:precision} shows that any such escape requires
    the oscillation $\max g_n - \min g_n$ to grow without bound in $n$.
  \item \emph{Non-residue local features.}  Potentials of the form
    $\log_2 n + h(t(n))$ where $h \colon \N \to \R$ depends on the
    trailing-$1$ count of $n$ are not directly covered, although a
    similar argument can be carried out: along the family
    $n_M = 2^{\ell + k} M - 1$ both $t(n_M)$ and $t(S^\ell(n_M))$ are
    pinned by construction, so $h$ contributes a fixed amount on each
    side of~\eqref{eq:diff-V} and the same lower bound applies.  We
    leave the full statement to the interested reader.
  \item \emph{Block structure and ordinal potentials.}  Goodstein-style
    ordinal potentials~\cite{KirbyParis1982} or more elaborate
    block-structure potentials are not covered.  Indeed, the Goodstein
    function admits a successful ordinal potential (with values in
    $\varepsilon_0$) showing termination, even though the function is
    Peano-arithmetic-unprovably-terminating.  Whether an analogous
    transfinite potential exists for Collatz is open.
  \item \emph{Non-local potentials.}  Anything depending on the entire
    history of an orbit, or on global features of $n$, escapes our
    framework.
\end{enumerate}

\paragraph{Relation to Conway's universality.}
Conway~\cite{Conway1972} proved that the natural class of
\emph{generalized} Collatz maps (piecewise linear with rational
coefficients on residue classes) is computationally universal; in
particular, no algorithmic procedure can decide convergence for all
such maps.  Our Theorem~\ref{thm:main} is in the same spirit: just as
Conway rules out a class of \emph{decision procedures}, we rule out a
class of \emph{Lyapunov certificates}.  The two results are
complementary: Conway shows generalized Collatz is universal; we show
that one specific certificate scheme is intrinsically too weak even
for the standard $3x+1$ map.

\paragraph{Relation to existing potential-function attempts.}
Many authors have explored small-$k$ instances of $V_{k, g}$ and
observed failure on specific residue classes (e.g., $n \equiv 7
\pmod 8$, $n \equiv 15 \pmod{16}$); the present paper consolidates
these observations into a uniform statement and a single conceptual
obstruction.  We do not claim that the impossibility for $k = 2, 3$
is new; we claim that the proof for \emph{all} $k$ is, and that the
identification of the obstruction with the trailing-$1$ blocks
clarifies why ad-hoc attempts at low $k$ inevitably failed.

\paragraph{Open problem.}
What is the smallest growth rate $f(k)$ such that some choice of
$g_k \colon \Z / 2^k \Z \to \R$ with $\max g_k - \min g_k \leq f(k)$
becomes a strict Lyapunov function for $S$?
Corollary~\ref{cor:precision} shows that bounded $f$ cannot work, but
gives no upper bound on the necessary growth.  A quantitative
statement (e.g., $f(k) = \Omega(\log k)$ or $f(k) = \Omega(k)$) would
give a sharper picture of how non-local any successful Lyapunov
function must be.

\paragraph{Reproducibility.}
Theorem~\ref{thm:main} can be verified numerically by choosing any
$k$ and $g$, then computing $V_{k, g}(S(n_M)) - V_{k, g}(n_M)$ along
the trailing-$1$ family $n_M = 2^{\ell + k} M - 1$ for $M = 1, 3, 5,
\dots$ and $\ell = \lceil (G + 1) / \log_2(3/2) \rceil$.  All
arithmetic is exact; no numerical issues arise for $k \leq 20$ and
$M \leq 2^{30}$.

%% =================================================================
\section{Conclusion}
%% =================================================================

We have shown that no finite local potential function of the form
$\log_2 n + g(n \bmod 2^k)$ with $k$ fixed can serve as a strict
Lyapunov function for the accelerated Collatz map $S$.  The obstruction
is direct: the trailing-$1$ block of length $\ell$ produces growth
$\ell \log_2(3/2)$ in $\log_2 n$, while the correction $g$ contributes
at most a fixed constant.  As a consequence, any successful Lyapunov
attack on the Collatz conjecture must employ either adaptive precision
($k = k(n) \to \infty$), non-residue features (e.g., trailing-$1$
counts), or genuinely non-local information (e.g., orbit history).
Our result delineates the boundary of what the most natural local
ansatz can achieve, leaving open whether any of the more elaborate
variants succeeds.

%% =================================================================
\begin{thebibliography}{99}
%% =================================================================

\bibitem{Conway1972}
J.\,H.~Conway,
\emph{Unpredictable iterations}, in:
\emph{Proc.\ Number Theory Conf.\ (Boulder, 1972)}, pp.\ 49--52,
Univ.\ Colorado, Boulder, CO, 1972.

\bibitem{KirbyParis1982}
L.~Kirby and J.~Paris,
\emph{Accessible independence results for Peano arithmetic},
Bull.\ London Math.\ Soc.\ \textbf{14} (1982), 285--293.

\bibitem{Lagarias1985}
J.\,C.~Lagarias,
\emph{The $3x+1$ problem and its generalizations},
Amer.\ Math.\ Monthly \textbf{92} (1985), 3--23.

\bibitem{Lagarias2010}
J.\,C.~Lagarias (ed.),
\emph{The Ultimate Challenge: The $3x+1$ Problem},
American Mathematical Society, Providence, RI, 2010.

\bibitem{Tao2022}
T.~Tao,
\emph{Almost all orbits of the Collatz map attain almost bounded values},
Forum Math.\ Pi \textbf{10} (2022), e12.

\end{thebibliography}

\end{document}
